\chapter{Related Work}
\label{ch:related}

The centralized platforms of \Cref{sec:traditional-pm} are mature and widely used, so one may ask whether software project management still needs new tools. This chapter argues that it does for one specific guarantee: an auditable, tamper-evident history of project and task changes that does not depend on trusting a single provider. It reviews the prior work that addresses this guarantee with blockchain technology and uses it to position the contribution of the thesis.

% ─────────────────────────────────────────────
\section{Decentralized Data Management with Hyperledger Fabric and IPFS}
\label{sec:adel}

The work closest to this thesis is that of \textcite{adel2023}, who propose a decentralized data management system for the construction industry built on Hyperledger Fabric and IPFS. Their setting has the same structure as the one addressed here: a construction project involves several independent parties (contractors, architects, engineers) who share task metadata and documents and need a verifiable, tamper-evident record of changes without a central authority.

\textcite{adel2023} separate the data in the same way as this thesis. Hyperledger Fabric stores structured metadata and the task and progress records on-chain, which gives an immutable history of updates, and large binary files such as drawings, images and documents are stored on IPFS and referenced from the ledger by their content identifiers (CIDs). Large files therefore stay out of the blockchain state, and every file reference is still bound cryptographically to the exact content of the file. They also use Fabric's channels and private data collections to restrict some information to subsets of participants, and the ledger state changes only through chaincode functions that handle project updates, task assignments and asset tracking.

This work matters for the thesis for three reasons. It shows that the Fabric-plus-IPFS architecture has already been validated for decentralized project data in practice. It establishes the convention of on-chain metadata and off-chain files that this thesis follows for task attachments (\Cref{ch:method}). And it targets construction, where the task model (physical progress milestones, drawings, regulatory documents) is very different from the software task model addressed here (issues, priorities, status transitions in a todo/in-progress/done workflow, artifacts close to the code). The architecture of this thesis is similar to that of \citeauthor{adel2023}, but the task model is different, and the target users are software teams used to tools such as Jira or GitHub Projects. Both differences place their own requirements on the data model of the chaincode and on the workflow presented to users (\Cref{ch:method}).

% ─────────────────────────────────────────────
\section{Blockchain-Based Project Management}
\label{sec:related-pm}

Several reviews have surveyed the use of blockchain in project management. \textcite{sonmez2021review} review the use cases proposed in the literature and derive a decision framework for designing such systems. \textcite{noteboom2026coordination} analyze 29 studies and find that adoption is strongest where coordination between several parties and verification by several parties matter most, as in integration, communication, quality and stakeholder management, and that outside construction it is still mostly conceptual or experimental. \textcite{seru2025design} reviews 58 publications and finds that nearly 70\% of them are design-science studies that build a prototype, while fewer than 10\% rely on an explicit theory. This thesis belongs to the first group, since it builds an artifact and evaluates it.

\paragraph{Systems on Hyperledger Fabric.} The closest system in terms of platform is the master's thesis of \textcite{batista2021}, a proof-of-concept project management tool on Hyperledger Fabric. Its aims are to give clients and other external stakeholders the same ``single source of truth'' as managers and developers, and to make audits easier. Its Node.js chaincode is organized into five smart contracts for projects, users, tasks with dependencies, timesheets approved by managers, and an automated risk analysis that detects late projects and tasks, budgets close to their limit, and developers with too many or too few tasks. Every user receives a Fabric identity from the certificate authority, and the system uses CouchDB to query records by field value. The work provides an API but no user interface, and its evaluation measures performance only (\Cref{sec:related-performance}). Fabric-GC~\cite{li2022fabricgc} is a Gantt chart system on a consortium Fabric network: several organizations plan a shared project, and access control decides who can see which part of the plan.

\paragraph{Agile and Scrum on blockchain.} ChainAgile~\cite{qureshi2024chainagile} is aimed at distributed Scrum teams. Smart contracts on a private Ethereum network manage the sprint backlog, user stories, requirement verification, acceptance tests and payments, with automatic penalties for late payments and missed deadlines, and IPFS provides off-chain storage. Scrum-Chain~\cite{silva_scrumchain} records the events of the Scrum cycle (backlog, team, sprints and tasks) through four Solidity contracts on an Ethereum test network. It writes every operation both to a PostgreSQL database and to the blockchain, and users sign transactions with the MetaMask wallet. It is one of the few such systems evaluated with users: ten student teams used it for a practical project over several sprints. In the questionnaires, 75\% rated the experience as excellent and 25\% as good, and the audit trail was named as its most useful feature. The main problems reported were the wallet integration, server errors and the number of steps needed for simple actions.

Other works address narrower problems. \textcite{gupta2021allocation} use smart contracts to allocate tasks to the members of distributed software teams through bidding. \textcite{rocky2021requirements} propose a hierarchical permissioned blockchain that controls which stakeholders may change which requirements in an agile project and keeps a traceable record of every change. Closest to the guarantee pursued in this thesis, \textcite{demi2024traceability} propose a blockchain-based knowledge base that records what was created or changed in a software project, how, when and by whom, and collect feedback on it from experts.

\paragraph{Relation to this thesis.} A tamper-evident record of project activity is a goal shared by all these works, but each differs from ChainBoard in at least one important way. Scrum-Chain requires users to hold wallets and pay transaction fees, and its own evaluation names the wallet as the main usability problem; ChainBoard runs on a permissioned network where users need neither. Scrum-Chain also keeps a second copy of the data in a conventional database, which the interface reads, while in ChainBoard the ledger is the only store. The Fabric-based tool of \textcite{batista2021} has no user interface and gives users no way to check a history, while ChainBoard has a board interface, file attachments on IPFS and an audit verification whose result can be re-checked offline. Some of these systems also go further than ChainBoard: Batista gives every person their own Fabric identity, and ChainAgile and Scrum-Chain support the full Scrum framework. Both points are discussed among the limitations of this thesis (\Cref{sec:limitations}).

% ─────────────────────────────────────────────
\section{Other Blockchain-Based Data Management Systems}
\label{sec:other-blockchain}

Blockchain-based data management has also been explored outside construction, for example for supply chain traceability and healthcare records. These domains have the same basic requirement as the work of \textcite{adel2023} and this thesis: several parties who do not trust each other need a shared, verifiable history of state changes. Such systems usually prefer permissioned blockchains, for the reasons given in \Cref{sec:blockchain}: the participants are known in advance, the required throughput is moderate, and organizations often need to keep some data private from each other. Their data models are quite different from project management (they may track the custody chain of a physical product, for instance, instead of the lifecycle of a software task), but they use the same pattern of anchoring off-chain content on an immutable ledger through content-addressed identifiers. They therefore support the architectural choice described in \Cref{sec:adel} and adopted in \Cref{ch:method}.

The Fabric-plus-IPFS pattern also appears in recent work on document management. \textcite{wen2024sharing} store data on IPFS and record its content identifier, its metadata and a log of data operations on a Fabric ledger, and use Fabric's private data to share data between groups of users. \textcite{calderon2026gestion} build a prototype for the long-term preservation of electronic documents on Fabric and an IPFS cluster, and report that Fabric's transaction times stay low and constant as the chain grows. In construction, \textcite{das2022document} propose a blockchain-based document management framework, which \textcite{adel2023} discuss among their prior work. \textcite{hushare2025practices} draw software engineering recommendations from building a Fabric prototype for supply chain traceability. None of these works targets software project management, but together they show that anchoring off-chain files to a permissioned ledger is an established design.

% ─────────────────────────────────────────────
\section{Tamper-Evident Logging}
\label{sec:related-logging}

Tamper-evident logging has been studied independently of blockchains (\Cref{sec:hash-chains}). \textcite{schneier1999secure} and \textcite{crosby2009tamper} study logs whose integrity can be checked after an attack, and Certificate Transparency~\cite{rfc6962} applies the approach at internet scale. These systems protect a \emph{log}, usually a single append-only sequence kept by one operator and checked by auditors. ChainBoard differs in two ways. Its history is kept by a replicated ledger that several organizations operate under consensus, and there is a separate history \emph{for each task and project}, tied to the object the user is looking at. The digest chain of \Cref{sec:audit-verification} takes the hash-chain idea from this literature and applies it to a single object, so that a downloaded copy of its history can be checked offline.

Blockchains have also been used directly to store audit logs. BlockAudit~\cite{ahmad2018blockaudit} turns the entries of a conventional audit log into transactions on Hyperledger and measures the cost in latency, network size and payload size. \textcite{belchior2020audits} store audit logs on a permissioned consortium blockchain and run automated audits over them. These systems protect a general-purpose log. In ChainBoard the ledger is the primary store of the project data itself, so the audit trail cannot drift apart from the data it describes.

% ─────────────────────────────────────────────
\section{Performance of Hyperledger Fabric}
\label{sec:related-performance}

The performance of Fabric has been studied in its own right. \textcite{thakkar2018performance} benchmark the platform and show that throughput and latency depend strongly on the block size and timeout of the ordering service, the endorsement policy, the number of concurrent clients, and conflicts between transactions that touch the same keys. These results shaped the data model of this thesis, which gives every account, project and task its own key so that unrelated writes do not conflict. They also determine how its measurements should be read (\Cref{sec:results-performance}): the figures describe one machine under a synthetic load, and show orders of magnitude and the behavior under contention, not the limits of the platform.

The project management tool of \textcite{batista2021} gives a point of comparison for a similar workload. In load tests through its API on a single laptop, all read requests succeeded up to 400 concurrent requests, while the success rate of writes fell from 100\% at 200 concurrent requests to 65\% at 400. A direct benchmark of the chaincode with Hyperledger Caliper reached at most 41.6 write and 159.1 read transactions per second, with read latencies around 60\,ms and write latencies between 3.5 and 12.4\,s. Batista attributes most of the slowdown to the communication between client and network and to the hardware, and concludes that the tool is adequate for a small or medium-sized company.

% ─────────────────────────────────────────────
\section{Positioning of This Thesis}
\label{sec:positioning}

Traditional platforms such as Trello, Jira, Asana, GitHub Projects and GitLab (\Cref{sec:traditional-pm}) offer mature project management functions but no cryptographic guarantee that the history is intact: where audit logs exist, the provider controls them and could in principle alter them without anyone being able to check. \textcite{adel2023} show that a Fabric-plus-IPFS architecture can close this gap, but in a specific domain whose data model does not carry over to software project management. Of the blockchain-based project management systems reviewed in \Cref{sec:related-pm}, those on Ethereum depend on the Ethereum model of accounts and transactions, which in Scrum-Chain means wallets and fees~\cite{qureshi2024chainagile,silva_scrumchain}; Scrum-Chain also keeps a second copy of the data outside the ledger~\cite{silva_scrumchain}; and the Fabric-based tool of \textcite{batista2021} stops at an API, without an interface or a way for users to check a history.

This thesis adapts the validated Fabric-plus-IPFS pattern to software project management. Projects, tasks, role-based membership and status transitions are chaincode-governed state, and the system keeps the functions that users of Trello or Jira expect: task boards, assignments, priorities and file attachments. Beyond the data model, it adds three things not covered by the work above: permissions enforced by the chaincode itself according to member roles, user accounts recorded as ledger transactions, and an audit verification whose result can be re-checked offline, in the spirit of tamper-evident logging (\Cref{sec:related-logging}). It also states where these guarantees end (\Cref{sec:results-threats}). \Cref{ch:method} describes the design and implementation.
